% =====================================================================
%  THE LEDGER FRAMEWORK — Valuation Manifesto, Version 1.0
%  Governing principles for valuation and profit-and-loss.
%  Subordinate to the Architectural Constitution
%  (ledger_manifesto_v1_4.tex) and the Market Data Manifesto 1.1
%  (MarketDataManifesto_1.1.tex). Rated against both; where it
%  conflicts with either parent, the parent prevails and the conflict
%  is parked with exact text.
%  Drafting lead: THORP. Rigor-check: FORMALIS. Binding support
%  memos: JACOBI (the PnL calculus), GATHERAL (the inheritance
%  interface). Review: KLEPPMANN, TALEB. Certification: CONCORDIA.
%  Source material (not a constraint): the Quantitative Volatility
%  Corpus, Volumes I--II (Valuation/volume_I.tex, volume_II.tex).
%  Build: pdflatex (twice). Standard TeX Live; lmodern and microtype
%  used when available and skipped otherwise.
% =====================================================================
\documentclass[11pt,a4paper]{article}

\IfFileExists{lmodern.sty}{%
  \usepackage[T1]{fontenc}\usepackage{lmodern}%
  \IfFileExists{microtype.sty}{\usepackage{microtype}}{}%
}{}
\usepackage[margin=2.5cm,marginparwidth=2.0cm,marginparsep=2mm]{geometry}
\usepackage{parskip}
\usepackage{amsmath}
\usepackage{amssymb}
\usepackage{booktabs}
\usepackage[hidelinks]{hyperref}

% --- article addressing: margin identifiers, mirroring the parents ---
\IfFileExists{xcolor.sty}{\usepackage{xcolor}}{}
\ifdefined\textcolor\else\def\textcolor#1#2{#2}\fi
\newcommand{\vlabel}[1]{\leavevmode\hypertarget{vm-#1}{}%
  \marginpar{\raggedright\scriptsize\ttfamily\textcolor{gray}{#1}}}
\newcommand{\pelabel}[1]{\leavevmode\hypertarget{pe-#1}{}%
  \marginpar{\raggedright\scriptsize\ttfamily\textcolor{gray}{#1}}}

\setcounter{secnumdepth}{0}
\providecommand{\tightlist}{\setlength{\itemsep}{2pt}\setlength{\parskip}{0pt}}

% traceability tag at the close of each article
\newcommand{\trace}[1]{\par\smallskip{\small\itshape (#1)}}

% --- Part B (The Pricing Engine) notation, reused from the corpus by name ---
\newcommand{\dd}{\mathrm{d}}
\newcommand{\Vanna}{\mathrm{Vanna}}
\newcommand{\Volga}{\mathrm{Volga}}
\newcommand{\Gadj}{\Gamma_{\mathrm{adj}}}
\newcommand{\smod}{\sigma_{\mathrm{mod}}}
\newcommand{\sprod}{\sigma_{\mathrm{prod}}}
\newcommand{\Dhyb}{\Delta_{\mathrm{hyb}}}
\newcommand{\Ghyb}{\Gamma_{\mathrm{hyb}}}
\newcommand{\satm}{\sigma_{\mathrm{ATM}}}
\newcommand{\ssm}{\sigma_{\mathrm{smile}}}
\newcommand{\lF}{\ell_F}
\newcommand{\Fref}{F_{\mathrm{ref}}}
\newcommand{\volslope}{\mathrm{vol.slope}}

\begin{document}

\begin{center}
  {\LARGE\bfseries THE LEDGER FRAMEWORK}\\[0.9em]
  {\large\bfseries Valuation Manifesto}\\[0.6em]
  {\normalsize Governing Principles for Valuation and Profit-and-Loss ---
   Subordinate to the Architectural Constitution and the Market Data
   Manifesto}\\[1.4em]
  Version 1.0 \quad---\quad July 2026
\end{center}

\vspace{0.6em}
\begin{center}\rule{0.35\linewidth}{0.4pt}\end{center}
\vspace{0.6em}

\begin{center}\itshape
A valuation is a projection of the record, and each mark earns its place
by proving itself from the one before.
\end{center}

\vspace{0.4em}

\begin{center}\textbf{Abstract}\end{center}
\begin{quote}
A valuation is not a bare number. It is the value of a unit read at
coordinates --- a ledger position, a unit state, a corporate-action
frame, a market-data cut, and the models the price is bound to --- and
the same coordinates reproduce it to the last minor unit. This manifesto
states the principles that make valuation and profit-and-loss a
projection of the one record and nothing beside it. Two layers, kept
apart on the first page, carry the whole design: each model-priced leg
\emph{re-enters as a recorded observation}, bound to its model by
lineage (MD-15); the net asset value, the profit-and-loss, and the
valuation chain that assembles them are a \emph{projection} of the record,
storing
nothing (C-8.2). A unit's current value is a \emph{valuation chain}: each
mark is proven from the last by a profit-and-loss explain whose residual
sits below a declared bound, and an unexplained move is a broken chain,
forbidden. Every mark is dispute-ready, settled by replay, extending the
Market Data Manifesto's exhibit list. A corporate action is a valuation
sandwich, its frame re-coordination a zero-profit identity and its
near-zero residual the proof it went well. A risk report is the same
valuation on shifted market data --- no second system, no second code
path. Each principle is an \emph{article} derived from the seven governing
principles the two parents already fixed --- auditability, reproducibility,
time travel, correctness, order, simulability, and the derived
dispute-readiness --- and rests on the clause of a parent that carries it.
The manifesto is in two parts. Part~A is the doctrine above --- the ten
articles that govern the record and its chain. Part~B, deliberately
distinct, steps beneath the doctrine to state and derive one mathematical
fact about the pricing function the chain stores: the price and the
sensitivities a book hedges with are drawn, in general, from two different
worlds, and the object that measures the gap --- the \emph{product
instantaneous volatility} --- is what the carry line of every explain must
reconcile against.
The document is rated against both parents, never the other way round;
where it conflicts with either, the parent prevails.
\end{quote}

\vspace{0.4em}
\begin{center}\rule{0.35\linewidth}{0.4pt}\end{center}
\vspace{0.6em}

\begin{center}
{\large\bfseries PART A \quad---\quad Valuation Doctrine}\\[0.3em]
{\small\itshape The record and its chain}
\end{center}

\section{1. Valuation in the Ledger's Terms}

A \textbf{valuation} is the value of one unit, together with its
sensitivities, read at a fixed set of coordinates. It is never the truth
about what the unit is worth; it is what a stated model says the unit is
worth given stated inputs, exactly as the Market Data Manifesto holds an
observation to be a fact about what a source said, not about what a thing
is (MD-1). The value of the whole book is a sum of such per-unit
valuations over the owned coordinate (C-4.10, C-8.2).

\textbf{Two layers, and they must never be confused.} A stored per-unit
model price, or a stored greek, is a \textbf{re-entered observation}: the
ledger runs no model (C-14.9), so a model is evaluated outside the fold
and its output re-enters through the one door as a recorded observation
(C-14.15), bound to that model by lineage (MD-15), and stale the moment an
input it consumed moves (MD-8). The net asset value, the profit-and-loss,
and the \textbf{valuation chain} that assembles them are of the other
kind: a \textbf{projection} that computes over the record --- owned
quantity times price, summed over valued units --- storing nothing and
rebuilding from the record alone (C-8.2, C-14.11). The projection consumes
the model-priced legs as \emph{leaves} (MD-12), exactly as it consumes a
directly observed stock quote: a plain observation enters the same sum. So
``valuation is a pure function of ledger state'' is a true statement about
the projection \emph{recipe}; ``a valuation re-enters as an observation''
is a true statement about each model-priced \emph{leg}. Both hold at once,
at their own layers, and every principle below follows from the pair. The
distinction is not pedantry; it is where a defect is repaired. A wrong leg
is a mismarked observation, re-derived with its model; a wrong total is a
mis-assembled projection, rebuilt from the record. Confuse the two and you
repair the wrong thing.

\textbf{Vocabulary, mapped once.} This manifesto uses the fixed names of
both parents and coins no synonym for any of them (C-Auth.4). It reuses
\emph{unit}, \emph{wallet}, \emph{balance}, \emph{projection}, \emph{the
immutable log}, \emph{the market data operator}, \emph{profit-and-loss
explain}, \emph{explain item}, and NAV; \emph{observation}, \emph{re-entered
observation}, \emph{frame}, \emph{cut}, \emph{as-of}/\emph{as-at},
\emph{complete lineage}, and \emph{simulated path} come from the Market
Data Manifesto unchanged. It coins four names, each checked against the
fixed set: the \textbf{valuation record}, the \textbf{valuation chain}
(always two words --- the immutable log alone is hash-chained, C-4.8, and
``chain'' bare is reserved for it), the \textbf{profit-and-loss explain
certificate}, and the \textbf{corporate-action valuation sandwich}. A
scenario perturbation is a \textbf{shift} --- never a ``shift operator'',
since \emph{the market data operator} names one thing only, the
corporate-action transform of C-9.2 (MD-13). A scenario is a
\emph{simulated path} under a recorded shift (MD-11); no ``derived world''
is coined.

\section{2. The Articles}

Each article derives from one or more of the seven governing principles
and rests on a clause of a parent --- the Constitution (C-\emph{sec}.\emph{n})
or the Market Data Manifesto (MD-\emph{n}). The identifiers VM-\emph{n}
are append-only: never reused, never renumbered.

\vlabel{VM-1}\textbf{VM-1. A valuation is a re-entered observation of a
unit's price; the book value is a projection over such observations.} The
two layers of Section~1 are the first article, because every other rests
on them. Where a price is read directly from the record --- a listed close
--- it is a plain observation. Where a price is produced by a model, the
ledger does not run the model (C-14.9); the model runs outside the fold
and its output, together with the sensitivities computed with it,
re-enters through the one door as a recorded observation bound to that
model (C-14.15, MD-15) --- the \textbf{valuation record}. The net asset
value is then a projection: replay the log to recover the wallet's
composition, multiply by the current price vector over the valued units,
and store nothing (C-8.2, C-8.3). What the projection multiplies is a
mixture of plain and re-entered observations; it treats them alike, as
leaves (MD-12). No valuation adds a store to the ledger: it is a further
fold of the same log (C-1.4).
\trace{Correctness, Reproducibility; C-8.2, C-14.9, C-14.15, MD-6, MD-15.}

\vlabel{VM-2}\textbf{VM-2. A valuation is meaningless without its
coordinates; the same coordinates reproduce it bit for bit.} A number
alone is not a valuation. A valuation is fixed only by its coordinates: the
ledger position it values --- the composition at an \emph{as-of} world-moment,
recorded at an \emph{as-at} cut (C-2.7, MD-4); the lifecycle state of the
unit, which pricing depends on and which is read from the ledger, never
from a store of its own (C-8.3); the corporate-action \textbf{frame} the
value lives in (MD-13); the market-data \textbf{cut} that fixes which
observations are in force (MD-6, MD-12); and the models the price is bound
to (MD-15). The position's as-at cut and the market-data cut are two
coordinates, not one, and independently selectable: hold the position at
yesterday's as-at and move the market-data cut to today, and you have
yesterday's book on today's prices --- a distinct, legal query on which
profit-and-loss attribution and the shifted worlds of VM-10 both rely.
Given the same coordinates, any party recomputes the identical
valuation, to the last minor unit (C-2.2). Reading a recorded valuation
back is unconditional; \emph{re-deriving} it from its inputs needs the
retained model --- the one coordinate whose retention is external to the
record, a governance matter (C-14.15, \S4). Because each input is pinned at
its own cut and never to a mutable store, a valuation struck ``mid-update''
cannot arise (MD-8, MD-12): the state a valuation reads is the folded
state, whole, or the valuation is not struck. A valuation carried without
its coordinates is not a weaker valuation --- it is not one.
\trace{Reproducibility, Time travel, Correctness; C-2.2, C-8.3, MD-4, MD-13.}

\vlabel{VM-3}\textbf{VM-3. A unit's value is a valuation chain; each mark is
proven from the last.} The current valuation of a unit is not an isolated number
but the head of a \textbf{valuation chain} --- an ordered sequence of
valuation records, each proven from the one before it by a
profit-and-loss explain (VM-4). The chain is append-only: a re-mark is a
new link --- itself a re-entered observation (VM-1) --- forward, never an
edit of an old one (C-4.8, C-12.4), and the
value the book showed at any past cut survives unrewritten beside the
current one (C-12.1). The greeks that close one link seed the next: today's
exit sensitivities are tomorrow's entry sensitivities, so the explain
carries measured sensitivities at both ends of every link. Measured
endpoint greeks are not a convenience: they keep the first two orders
exact and squeeze any gap in Hessian coverage --- the cross-greeks an
implementation did not populate --- into the third-order term and the
residual, rather than smearing it through delta and vega, where it would
corrupt the very hedge the greeks are supposed to size. The chain is
the valuation layer's instance of complete lineage (MD-6); it reaches back
through the price vector and the operator-adjusted inputs to the leaf
observations. A unit may carry more than one valuation chain at once ---
one per re-marking convention, cadence, or desk --- and they do not
compete: each decomposes the one change in net asset value (C-8.4) and
localises to its own recorded convention coordinate (VM-5).
\trace{Order, Auditability; C-4.8, C-12.1, C-12.4, MD-6.}

\vlabel{VM-4}\textbf{VM-4. The profit-and-loss explain is the certificate,
and it states honestly what is exact.} Each link carries a
\textbf{profit-and-loss explain certificate}: a decomposition of the value
change between two marks, read the way a desk reads it. Carry first --- the
value the position sheds or earns as a day passes at unchanged market.
Then the delta line --- the move read against the entry delta. Then the
gamma-and-theta balance, the core of the daily explain: convexity earns on
the move, decay pays for it, and a delta-hedged position turns a profit
only once the move clears its breakeven, the point where the day's gamma
gain just covers the day's theta. That balance closes at exactly one
volatility --- the \emph{product instantaneous volatility} of Part~B (PE-4).
The carry line therefore carries a declared, recorded term naming \emph{which}
of the volatilities in play it was reconciled against --- the model's, the raw
implied-surface, or the product instantaneous volatility. That term is a
governance lever --- it decides whether the Feynman--Kac discrepancy (PE-2) sits
in the carry line or lands in the residual --- so a change to it is an auditable,
reviewable event, exactly as the attribution convention (VM-5) and the
certification bound (VM-6) are. Then vega, bucket by bucket --- by name,
by expiry, by strike, never a single number --- with the cross-greeks
vanna and volga among the higher-order terms. Then the lifecycle and
state changes the link crossed. And last, one explicit residual (VM-6).
Two forms of the market lines exist and the certificate must not conflate
them. The \emph{ex-ante} estimate --- the reading available \emph{before}
the move, computed from the entry greeks alone --- is
$\Delta\!\cdot\!\delta + \tfrac12\,\delta^{\mathsf T}\Gamma\,\delta$, the
delta-and-gamma a desk pre-computes, and carries a
\emph{deterministic, signed} gap to the reprice equal to the variation of
gamma over the move. The \emph{exact} decomposition uses the greeks at
\emph{both} ends: the value change is the move times the average of the
entry and exit deltas, less a one-twelfth correction for the gamma that
moved (a both-endpoint quadrature correction), plus a residual $R$ that is
identically zero only for a book cubic
in its risk factors. Using entry greeks in place of both-end greeks ``is
not a smaller model; it is an inconsistent one'', so the certificate names
which form it used. The residual behaves well only under conditions: for a
value smooth to fourth order, $R = O(\lVert\delta\rVert^4)$ --- doubling
the move multiplies its leading part by sixteen; at a digital, a barrier,
an autocallable, any payoff that is not smooth, $R = O(1)$, the size of
the discontinuity, the instant a move straddles the feature; and the
calendar direction is never covered by that bound --- a day of carry near
expiry lands in $R$ at full size. The certificate therefore \emph{never}
promises a universally bounded residual. The total the certificate explains is path-independent: it is exactly the
change in net asset value (C-8.4). Only the \emph{decomposition} depends on
the marks chosen along the way and on the convention (VM-5) --- the identity
that the parts sum to $\Delta$NAV holds on every path; the split into parts
does not, and no second running total of realised profit-and-loss is kept
beside NAV to disagree with it (C-8.4). Because the residual is the
balancing term --- the value change less every attributed line --- the
certificate reconciles to the change in net asset value by construction,
not by coincidence: it is a projection of that one number, and no split it
prints can make it disagree with the fold. Attribution decomposes
$\Delta$NAV; it never recomputes it. This explain \emph{composes} with
the Constitution's reordering explain: there is one profit-and-loss
explain, and a reordering-refold line (C-12.6) and a market-factor line
here are two disjoint kinds under one roof, each attributed to exactly one
cause. The certificate extends the taxonomy; it does not re-coin the term.
\trace{Correctness, Auditability, Order; C-8.4, C-12.6, MD-5.}

\vlabel{VM-5}\textbf{VM-5. Attribution is convention-dependent; the
convention is a declared, recorded term.} How a move splits across the
greeks is not unique. Where a risk factor is not a market map of its
driver --- an implied basket correlation is the canonical case --- the
split between the spot channel and the vol channel depends on which
quantity the desk holds invariant as the others move, and two conventions
that reproduce today's marks exactly differ in tomorrow's attribution:
same marks today, different \emph{attributed} profit-and-loss tomorrow. That invariant is a
declared, recorded term of the valuation (MD-6). Two consequences follow.
The same convention reproduces the same attribution, bit for bit, a
deterministic function of the record. And a dispute about attribution is a
dispute about convention: replay \emph{localises} it to the differing term
(MD-14) but does not adjudicate which convention is right, which is model
choice and out of scope (MD-9, MD-15). The manifesto therefore never
speaks of \emph{the} attribution as unique, nor of a shadow greek --- the
value routed through the held-invariant factor --- as a property of the
market: under the engine default that holds the at-the-money volatility
constant, every spot-channel shadow number is identically zero, and must
be flagged as computed under that convention, never quoted as a market
fact.
\trace{Correctness, Reproducibility; MD-6, MD-9, MD-14, MD-15.}

\vlabel{VM-6}\textbf{VM-6. ``Explained'' means the residual sits below a
declared bound; a breach is a recorded fact, never absorbed.} A move is
\emph{explained} when its residual $R$ falls inside a declared tolerance.
The tolerance is a declared term, not a universal constant: a valuation is
certified only when each instrument's residual is within its bound and the
aggregate is within the aggregate bound, and those bounds are recorded in
the model configuration, mapping between price space and vol space through
vega. These bounds roll up: the per-unit residuals aggregate to a
book-level residual against the aggregate bound, where cross-unit and
correlation profit-and-loss --- routed under the declared convention
(VM-5) --- must also come to rest. The residual is itself a re-entered
observation --- a recorded
diagnostic, stale if its inputs later move, never a silent pass (MD-9,
MD-15). One structural source of it is named in Part~B: a carry line that
books the model theta against the gamma the book carries but reconciles it at a
volatility other than the product instantaneous volatility (PE-4) leaves the
Feynman--Kac discrepancy (PE-2) sitting in $R$. Which volatility reconciles the
line is a declared choice (VM-4); making it silently --- mixing the two worlds
without declaring --- is exactly the unexplained profit-and-loss this article
refuses to pass. It is a sensor, not a rounding error: a residual that is small,
mean-zero, and regime-independent certifies the bookkeeping is a control
variate --- a correction whose smallness can be trusted ---; a residual
that is structurally large, sign-biased, or wakes up
in stress reports that the honest attribution of this position is a
scenario, not a greek, and demands re-mark or model review. Two things the
bound must survive. First, the anti-plug: that the certificate reconciles
to the change in net asset value proves \emph{nothing} --- the residual is
the balancing term (VM-4), so the parts sum to $\Delta$NAV whether the move
is explained or not; only the residual \emph{below its declared bound}
certifies. Reconciliation is automatic; explanation is earned. Second, the
bound's own adequacy is a governance assumption, named here as the parent
named the identifier grain and the delivery frame: a bound set too loose
makes every move ``explained,'' fires no breach, and empties ``unexplained
profit-and-loss is forbidden'' (VM-7) in silence. The guard is that the
residual is recorded whatever the bound --- its sign and regime behaviour,
not its level alone, are inspectable after the fact --- and that the bound's
calibration is itself a recorded, reviewable event, so a loosened bound is
an auditable change, never a silent one. A breached
bound is booked --- a named line in the profit-and-loss explain (MD-5) ---
and acted on. A bookkeeping that announces when to stop believing it is
worth more than one that is elegant while it lasts.
\trace{Correctness, Auditability; C-2.4, C-13.2, MD-9, MD-15.}

\vlabel{VM-7}\textbf{VM-7. Unexplained profit-and-loss is a broken chain,
and a broken chain is forbidden.} A link whose residual exceeds its
declared bound, or whose leaf inputs have moved under it, is not a link:
it is a broken valuation state, and a broken state cannot hide (MD-8). Two
mechanisms keep it visible. A projection over a corrected leaf recomputes
on the next read, so a stale \emph{projection} cannot arise. A re-entered
valuation record that consumed a corrected or superseded input --- a wrong
price, a restated corporate action behind its frame --- cannot re-run
itself (C-14.9); its lineage marks it stale and it stands as a flagged
open item for re-derivation (MD-8, MD-10), read through a projection that
carries the staleness flag, never as the raw stored number. Staleness
propagates forward along the chain: a superseded link flags stale every
downstream link whose certificate chained from it, and each re-proves
\emph{forward from the superseding predecessor} --- never editing the old
proof --- so a certificate can never keep attesting a transition from a
mark that no longer heads its segment (MD-6, MD-8). The as-known branch
survives beside the corrected one (C-12.1), so the chain resolves into the
two honest answers rather than silently forking. Nothing in the
chain is repaired by overwriting; a correction repairs forward, and where
money already moved it moves back only under authorised compensation
(C-12.4). What is forbidden is the \emph{silent} unexplained move: an
unexplained profit-and-loss is a defect the record must surface, not a
number the book may carry.
\trace{Correctness; C-1.5, C-14.9, MD-8, MD-10.}

\vlabel{VM-8}\textbf{VM-8. Every valuation is dispute-ready.} The current
valuation of every unit is dispute-ready: its chain, the certificates
along it, and its coordinates suffice to reproduce and defend it, bit for
bit, against any counterparty --- a valuation dispute and a collateral
dispute alike. This is the Market Data Manifesto's dispute-readiness
(MD-14) reached at the valuation layer, and it \emph{extends} that
manifesto's exhibit list: to the observation, its provenance, the cut, the
frame, and the bound model, valuation adds the \textbf{valuation chain}
and its \textbf{certificates}, exhibited as further rows of the one
replay, never a second replay doctrine. Its bound is exact and inherited:
replay settles whether the mark is the faithful, deterministic consequence
of the inputs the record held at its stated coordinates, the as-at pinned
to the mark's own cut so the number that stood then is reproduced, not a
later one. It does not settle a two-sided economic dispute --- a
counterparty's mark, from its own surface and convention, replays bit for
bit too --- but exhibiting the chain and the convention \emph{localises}
the disagreement to the differing input, frame, model, or convention,
which is the value replay delivers (VM-5, MD-14). A collateral dispute is
bounded the same way: replay settles the \emph{mark} inside the call, never
the collateral-agreement terms --- haircut, eligibility, netting --- which
are model choice and out of scope (C-Scope.11).
\trace{Dispute-readiness (Auditability + Reproducibility); C-2.1, C-2.2,
  MD-14, MD-15.}

\vlabel{VM-9}\textbf{VM-9. A corporate action is a valuation sandwich; its
frame re-coordination is not profit-and-loss.} Every corporate action that
changes a unit's state triggers a \textbf{valuation sandwich}: a valuation
struck immediately before the transition and another immediately after,
with a full profit-and-loss explain across the pair. Across the sandwich
the state and the coordinate frame change discontinuously, so the smooth
explain of VM-4 does not apply --- there is no smooth move to expand along.
The only admissible explanation of the jump is the market data operator
(C-9.2, MD-13): it transports the pre-frame mark into the post-frame,
acting on leaf observations, with derived objects recomputed from
operator-adjusted inputs, never scalar-transported. The frame
re-coordination itself is a \emph{zero-profit identity}: the same value in
new coordinates, the consistency-of-reference that forbids multiplying a
new share count by a stale price (C-11.3). Only genuine economics leaving
or entering the position --- a special dividend paid away --- is
profit-and-loss, and it earns its own explain line; there is no ``corporate
action profit-and-loss'' beyond it. That the sandwich residual is
$\approx 0$ is the proof the action went well, and $\approx$ means exactly
three declared tolerances and no more: the minor-unit rounding applied once
at read (C-4.6), the declared convention's own rounding, and the
calibration tolerance of any re-derived surface or curve (VM-6). A residual
materially above these is not noise --- it says the operator or the
delivery frame was mis-declared, which replay localises (MD-14) and MD-8
and MD-10 handle as a superseding event demanding re-derivation. The
operator exists only once the action's terms are \emph{resolved}, not
merely announced; before resolution the frame is provisional and the
sandwich is struck as such (MD-13). A third timing case is a resolved
action that reaches the door \emph{late}, after valuations were already
chained across its ex-date. This is a reordering, not a lineage
supersession: the intervening links never carried the action in their
lineage, so they are flagged stale by the reordering trigger (C-2.7,
MD-5), not by the frame-supersession path of VM-7. The action takes its
place at the ex-date, its sandwich is struck \emph{retroactively} there,
and the spurious market-move line the chain had booked --- a two-for-one
split reaching the door three days late reads as a fifty-percent crash
until it arrives --- is reclassified as the zero-profit frame
re-coordination, the reordering carried as its own line in the
profit-and-loss explain (C-12.6). The loss is never left standing.
\trace{Order, Correctness, Time travel; C-9.2, C-9.3, C-11.3, MD-13, MD-5.}

\vlabel{VM-10}\textbf{VM-10. Risk is valuation in simulated worlds.} A risk
report is not a separate system and not a second code path: it is the same
per-unit valuation recipe of VM-1 applied to \emph{shifted} market data.
The Constitution already fixes the doctrine --- ``we live in a simulation:
production is simply the one path whose events happened to be real''
(C-2.8) --- and the Market Data Manifesto that simulated market data is
real market data under a different seed (MD-11). A scenario is a simulated
path in which one or more observations are moved by a recorded
\textbf{shift}; a greek is a sensitivity read across the base and shifted
paths. Shifts \emph{compose} --- spot up twenty percent, volatility down,
in any combination --- each composition yielding a new vector of
valuations and greeks by re-application of the one recipe. Base and
simulated worlds share their units, their models, and the market data
operator; only the market-data state differs. And the gap between the
greek approximation and a full revaluation is itself a computable,
recordable fact --- the ex-ante estimation gap of VM-4: the deterministic,
signed gamma-variation term \emph{plus} the residual $R$ (VM-6), small for
a smooth book and $O(1)$ where the payoff is not --- so a sound system
computes both and records their difference, on the same terms as a
production day, never choosing silently between them. The shift is recorded, exactly as a seed
is (MD-11): it is the single non-record input of the scenario, so the path
replays. But a path that recorded only its shift could not replay --- every
model-priced valuation along it re-enters as an observation like any other
(VM-1), into the simulated path's \emph{own} record, never as a link in the
real unit's valuation chain --- so a stress run is a derived object with
full lineage, dispute-ready on the same terms as production (VM-8). And
because risk and the books run the one recipe, the risk report and the
official profit-and-loss cannot disagree by construction: the classic
risk-versus-books reconciliation break --- two systems, two code paths, two
numbers --- has no way to arise.
\trace{Simulability, Reproducibility; C-2.8, MD-11, C-14.11.}

\section{3. The Life of a Valuation}

One unit walks the whole picture. Follow a single long call, struck at
$100$ on a stock at $100$, three months to expiry, marked at a $30\%$
implied volatility. Its value is $5.98$, its delta $0.53$, its gamma
$0.0265$, its decay $-0.047$ per trading day. Each row of
Table~\ref{tab:chain} is a link; each is proven from the one above.

\begin{table}[h]
\centering
\small
\begin{tabular}{@{}lll@{}}
\toprule
\textbf{Link} & \textbf{What changed} & \textbf{Explain into this mark} \\
\midrule
Base & struck at all coordinates & value $5.98$; delta $0.53$, gamma $0.0265$ \\
Carry day & one trading day, spot flat & carry $-0.047$; residual $<0.001$ $\Rightarrow$ $5.93$ \\
Market move & spot $100 \to 102$, vol held & delta $+1.060$, gamma $+0.053$; \\
            &                              & ex-ante residual $-0.001$ (bound $0.010$) $\Rightarrow$ $7.09$ \\
CA sandwich & $2{:}1$ split; strike $\to 50$, & frame re-coordination (zero-profit); \\
            & deliverable doubles           & economics $0$; residual $\approx 0$ $\Rightarrow$ $7.09$ \\
\bottomrule
\end{tabular}
\caption{A valuation chain of four links on one unit. Each link isolates
one effect for legibility; on a real book carry and market move fall in a
single link and both appear as explain lines. The market-move residual is
the ex-ante gap (the higher-order and gamma-variation content), well inside
its declared bound. The corporate-action link is a sandwich: the pre-ex
value $7.09$ in the old frame equals the post-ex value $7.09$ in the new
one, and the near-zero residual is the proof the split was applied
consistently --- not a phantom valuation from a new share count against a
stale price (C-11.3).}
\label{tab:chain}
\end{table}

The market-move link also reads as the gamma-and-theta balance of VM-4:
the $\$2$ move clears this option's breakeven of $\$1.89$, so a
delta-hedged holder's gamma gain ($0.053$) just covers a day's decay
($0.047$). Inside the breakeven the day loses; outside it the day earns;
and that balance, not the delta, is what a long option position is really
holding.

\textbf{A shifted-world column.} From the base mark, a risk report is one
more valuation under a recorded shift. Shift spot $+20\%$, to $120$, and
re-apply the same recipe. The greek estimate is
$0.53\!\times\!20 + \tfrac12\!\times\!0.0265\!\times\!20^2 \approx +15.9$;
the full revaluation is $+14.9$ (the mark moves to $20.9$). The gap,
$\approx +1.0$ --- near a fifteenth of the move, because a twenty-percent
shift is far outside the regime where the quadratic is a good control ---
is not an error to be suppressed but a recorded diagnostic (VM-10). The
shift composes: add a volatility-down leg and the same recipe returns a new
vector, no second code path.

\textbf{A link that breaks.} Now swap the call for a digital that pays a
fixed coupon if the stock is above $100$ at expiry, and take a link that
straddles the strike a week out. The two-point explain fails on sight: the
reprice is the whole coupon, while the endpoint deltas, evaluated where the
value has already flattened either side of the strike, read near zero, so
the entire notional lands in the residual --- $O(1)$, not
$O(\lVert\delta\rVert^4)$. The certificate does not hide this: the residual
breaches its bound, the breach is booked as a named line (VM-6), and the
chain is flagged broken until the position is re-marked (VM-7). The honest
attribution here is a gap scenario with a probability attached, not a
greek --- which is exactly what VM-4 warned the smooth explain could not
see, made visible by the very residual it refuses to absorb.

\textbf{Disputed.} A counterparty contests the collateral mark. It is not
argued: replay exhibits the chain, its certificates, and its coordinates,
and reproduces the mark bit for bit, the as-at pinned to the mark's own cut
(VM-8, MD-14). Replay cannot decide whose surface or convention is right,
but it localises the disagreement to a differing input, frame, model, or
convention.

\section{4. What This Manifesto Does Not Govern}

This manifesto governs how valuation and profit-and-loss are recorded,
proven, and projected. It does not govern \textbf{which model, surface,
prior, or re-marking convention is correct} --- model choice, whose test is
hedging performance, out of scope (C-Scope.11, MD-9, MD-15); nor
\textbf{model governance, validation, and retention} --- a valuation
reproduces from recorded prices unconditionally and from a model's number
only if the model is kept (C-14.15); nor the \textbf{numerical methods} by
which a model is solved --- regression, quadrature, surrogate pricing ---
which are the reference implementation's, not the specification's (C-14.9);
nor \textbf{price formation}, the \textbf{terms of a corporate action}
(they arrive as recorded events, MD-13), or \textbf{settlement}. These lie
outside the Ledger boundary and are reconciled against it only at its
perimeter.

\section{5. Subordination}

This manifesto is subordinate to two parents --- the Architectural
Constitution and the Market Data Manifesto --- and rated against both,
never the other way round. Its vocabulary is theirs: one name per
component, and no synonym for any fixed term (C-Auth.4). Where a principle
here would conflict with either parent, the parent prevails and the
conflict is parked with the exact amendment text and the exact clause it
replaces --- never resolved by quiet narrowing, which for a guarantee of
either parent is the worse failure. A genuine departure amends this
document explicitly first; an article identifier VM-\emph{n} is never
reused or renumbered.

\section{Conclusion}

Valuation has the same root cause of failure as everything else the Ledger
touches: the same fact held in two places, marked by two logics, drifting
apart. The Ledger removes it here as everywhere. A valuation is a
projection of the one record; each model price it consumes is an
observation; each mark it produces is proven from the last by a
profit-and-loss explain whose residual it must keep below a declared
bound, and an unexplained move is a broken chain, not a number to carry.
Every mark carries its coordinates, is dispute-ready by replay, survives a
corporate action as a frame re-coordination that is not profit-and-loss,
and is re-run in a simulated world for risk by the very same recipe.
Operate it any other way and the divergence the architecture exists to
remove comes back; this way, a valuation is a fact, and there is nothing
left to reconcile.

\clearpage
\begin{center}\rule{0.35\linewidth}{0.4pt}\end{center}
\vspace{0.6em}
\begin{center}
{\large\bfseries PART B \quad---\quad The Pricing Engine}\\[0.3em]
{\small\itshape One fact about the pricing function the chain stores}
\end{center}
\vspace{0.4em}

Part~A governs the record and its chain: a mark is a re-entered observation
(VM-1), proven from the last by a profit-and-loss explain whose residual
sits below a declared bound (VM-4, VM-6). Part~B is \emph{deliberately
distinct}. It is not workflow doctrine; it steps beneath the doctrine to
state, and derive, one mathematical fact about the pricing function the
chain stores --- the fact beneath a single line of the explain, the line
that reconciles carry against convexity. That fact: the price and the
sensitivities a book hedges with are, in general, drawn from two different
mathematical worlds, and the object that measures the gap is a
\emph{product instantaneous volatility}. Part~B is descriptive of that
structure only. It proposes no model, no surface, no re-marking rule, and
no remedy beyond the single monitoring mandate of PE-6; the rule by which
an implied surface moves enters solely as a declared mathematical map,
never as an account of anyone's practice. Its articles are numbered
PE-\emph{n}, append-only, and cite Part~A by VM-\emph{n} and the corpus by
section and equation.

%% ---------------------------------------------------------------------------
\pelabel{PE-1}\textbf{PE-1. The objects: a pricing model, its Feynman--Kac
equation, and a declared surface-move map.} Fix a filtered probability space
carrying a risk-neutral measure $\mathbb{Q}$ and a strictly positive
underlying $S_t$. The following are standing for all of Part~B.

\textbf{(A1) Carry.} A constant short rate $r \ge 0$ and continuous dividend
yield $q \ge 0$; the drift of $S$ under $\mathbb{Q}$ is $(r-q)$. The corpus's
zero-rate specialization $r=q=0$, hence $F=S$, is flagged at each use.

\textbf{(A2) Pricing model $\mathfrak{M}$.} Its diffusion coefficient is a
function $\smod$: either \emph{constant} $\smod \equiv \sigma > 0$
(Black--Scholes for vanillas), or \emph{local} $\smod = \sigma_{\mathrm{loc}}(S,t)$
with $0 < \sigma_{\min} \le \sigma_{\mathrm{loc}} \le \sigma_{\max} < \infty$
and locally Lipschitz in $S$ uniformly in $t$ on compacts, so that
$\dd S_t = (r-q)S_t\,\dd t + \smod(S_t,t)\,S_t\,\dd W_t^{\mathbb Q}$ has a unique
strong solution.

\textbf{(A3) Product and price.} A product $\Phi$ with model price
$P(S,t) = \mathbb{E}^{\mathbb Q}\!\left[e^{-r(T-t)}\Phi \mid S_t = S\right]$,
assumed $C^{2,1}$ on the open domain
$\mathcal{O} = (0,\infty)\times[0,T)$. For a payoff with a discontinuity or a
kink in the state --- a digital, a barrier, an autocallable --- $\mathcal{O}$
excludes the discontinuity locus, and $C^{2,1}$ fails on it; this is the same
non-smoothness Part~A books as an $O(1)$ residual (VM-4), and PE-3 shows where
it re-enters here.

\textbf{(A4) Implied surface.} An arbitrage-free implied-volatility surface
$\Sigma_t$, calibrated to market at $t$ and valid in price space (MD-15); its
scalar value at the product's reference strike and expiry is the implied
volatility $\sigma$.

\textbf{(A5) The surface-move map (the hedging rule, as an object).} A declared
\emph{spot-dynamic} $\mathcal{D}$: a map that assigns to a spot level $S$ a
surface $\Sigma(\cdot\,;S)$, of class $C^2$ in $S$ near the valuation spot, so
that the reprice map
\begin{equation}\label{pe:reprice}
  g_{\Phi,\mathcal D}(S) \;:=\; P\bigl(S,\,\Sigma(\cdot\,;S)\bigr)
\end{equation}
is twice differentiable there. The moved surface stays admissible:
$\Sigma(\cdot\,;S)\in\mathcal{A}$, the set of arbitrage-free surfaces
(Vol~I, \S2.6), for $S$ in the neighbourhood; since $\mathcal{A}$ is open
(Vol~I, Def.~2.20), this holds automatically when the base surface is interior
to $\mathcal{A}$ and $\mathcal{D}$ is continuous into surface space.
$\mathcal{D}$ is a mathematical object only: sticky-strike, sticky-moneyness,
and sticky-delta are instances, as is the corpus dynamic --- the decomposition
$\sigma=\satm(\lF,T)+\ssm(m,T)$ (Vol~I, \S2.1.1, Eq.~(2.1)) \emph{held fixed}
as the forward $F=F(S)$ floats (Eqs.~(2.19)--(2.20), slope $\volslope$) --- each
a different $\mathcal{D}$. The decomposition (2.1) is a coordinate system, not a
dynamic; the dynamic is that system carried under the forward map. This map
$\mathcal{D}$ is not a new coordinate: it is the surface-dynamics content of the
re-marking convention Part~A already declares (VM-5) --- one declared, recorded
term (MD-6), read in two roles, the attribution split in Part~A and the hybrid
Greeks and $\sprod^2$ here. Accordingly the \emph{shadow gamma} of PE-5 is the
second-order instance of VM-5's \emph{shadow greek} --- the value routed through
the held-invariant factor --- not a separate object.

\textbf{Claim (Feynman--Kac).} Under (A1)--(A3), $P$ satisfies on $\mathcal{O}$
\begin{equation}\label{pe:fk}
  \partial_t P \;+\; \tfrac12\,\smod(S,t)^2 S^2\,\partial_{SS}P
  \;+\; (r-q)\,S\,\partial_S P \;-\; rP \;=\; 0,
  \qquad P(S,T)=\Phi .
\end{equation}
This is the Feynman--Kac PDE of $\mathfrak{M}$ --- the model's own pricing
equation, the consistency its price and Greeks satisfy by construction; its
diffusion coefficient is $\smod$, the volatility function of the model used to
price. Under (A2) the
standard Feynman--Kac hypotheses hold (uniform ellipticity and local
regularity of the coefficients), so \eqref{pe:fk} is exact, not a heuristic.
\trace{Correctness; VM-1, VM-2, MD-15; Vol~I \S1, Assumption~1.1.}

%% ---------------------------------------------------------------------------
\pelabel{PE-2}\textbf{PE-2. The discrepancy: hybrid Greeks do not satisfy the
model's Feynman--Kac equation.} The Greeks $\mathfrak{M}$ owns are the partials
$\partial_S P,\ \partial_{SS}P,\ \partial_t P$ of the same $P$; by construction
they combine to zero in \eqref{pe:fk}. The Greeks a book hedges with are, in
general, different objects: they are read off the reprice map
\eqref{pe:reprice}, in which the surface is moved by $\mathcal{D}$ and the
product is repriced.

\textbf{Definition (hybrid Greeks).} The \emph{hybrid delta} and \emph{hybrid
gamma} are the total derivatives of \eqref{pe:reprice} along the curve
$S\mapsto(S,\Sigma(\cdot\,;S))$,
\begin{align}
  \Dhyb &:= g_{\Phi,\mathcal D}'(S)
        = \partial_S P + \underbrace{D_\Sigma P[\partial_S\Sigma]}_{=:~\delta_\Delta},
        \label{pe:dhyb}\\[2pt]
  \Ghyb &:= g_{\Phi,\mathcal D}''(S)
        = \partial_{SS}P + \underbrace{2\,D_\Sigma(\partial_S P)[\partial_S\Sigma]
          + D_{\Sigma\Sigma}P[\partial_S\Sigma,\partial_S\Sigma]
          + D_\Sigma P[\partial_{SS}\Sigma]}_{=:~\delta_\Gamma},
        \label{pe:ghyb}
\end{align}
where $D_\Sigma P[\cdot]$ is the G\^ateaux derivative of $P$ in the surface
argument. These directional (G\^ateaux) derivatives are taken along the declared
$C^2$ curve $S\mapsto\Sigma(\cdot\,;S)$ of (A5); their existence is the standing
hypothesis that $P$ is twice differentiable in the surface argument there ---
realized concretely by the smooth Black--Scholes dependence of PE-5 --- not a
claim of Fr\'echet differentiability on surface space, which the manifesto
neither needs nor asserts. Coordinatized by strike,
$\delta_\Delta = \int (\partial P/\partial\sigma(K))\,(\partial\sigma(K)/\partial S)\,\dd K$
is the same object. The corrections $\delta_\Delta,\delta_\Gamma$ are the surface-response
terms; they are exactly the smile-adjustment terms of the total (routed)
Greeks of Vol~I, \S2.4.2 (Eq.~(2.21)), made explicit per instance in PE-5. Both
$\delta_\Delta$ and $\delta_\Gamma$ are proportional to the surface response
$\partial_S\Sigma$ that $\mathcal{D}$ declares.

\textbf{Claim.} With the model theta $\Theta := \partial_t P$ and the hybrid
Greeks $\Dhyb,\Ghyb$, the Feynman--Kac operator of \eqref{pe:fk} does not
vanish:
\begin{equation}\label{pe:residual-operator}
  \mathcal{L}_{\mathrm{hyb}}
  := \Theta + \tfrac12\,\smod^2 S^2\,\Ghyb + (r-q)S\,\Dhyb - rP
  \;=\; \tfrac12\,\smod^2 S^2\,\delta_\Gamma + (r-q)S\,\delta_\Delta,
\end{equation}
which is nonzero whenever the surface response is nonzero.

\textbf{Derivation.} Substitute \eqref{pe:dhyb}--\eqref{pe:ghyb} into
$\mathcal{L}_{\mathrm{hyb}}$ and add and subtract the model partials:
\[
\begin{aligned}
  \mathcal{L}_{\mathrm{hyb}}
  &= \underbrace{\bigl[\partial_t P + \tfrac12\smod^2S^2\partial_{SS}P
     + (r-q)S\,\partial_S P - rP\bigr]}_{=\,0\ \text{by \eqref{pe:fk}}}\\
  &\quad{}+\; \tfrac12\smod^2 S^2(\Ghyb-\partial_{SS}P)
  \;+\; (r-q)S(\Dhyb-\partial_S P).
\end{aligned}
\]
The bracket is the Feynman--Kac equation, identically zero; the remaining two
terms are $\tfrac12\smod^2S^2\,\delta_\Gamma + (r-q)S\,\delta_\Delta$, giving
\eqref{pe:residual-operator}. The discrepancy $\mathcal{L}_{\mathrm{hyb}} = 0$
exactly when the surface-response corrections $\delta_\Delta,\delta_\Gamma$
combine to zero. Sufficient are a surface-insensitive product
($D_\Sigma P = 0$, e.g.\ a forward) or a rule that holds the surface still
\emph{to second order} ($\partial_S\Sigma = \partial_{SS}\Sigma = 0$) --- note
$\delta_\Gamma$ retains the term $D_\Sigma P[\partial_{SS}\Sigma]$, so
$\partial_S\Sigma = 0$ alone does not suffice. The exact condition is that
$\mathcal{D}$ reproduce the model's own endogenous surface dynamics (Remark
below). $\square$

\textbf{Remark (two worlds, and when they coincide).} $P$ is the price under a
diffusion whose coefficient is $\smod$; in that world the implied surface is an
\emph{output}, and its response to spot is whatever the model dictates. The
hybrid Greeks encode an \emph{externally declared} response $\partial_S\Sigma$.
Equation \eqref{pe:residual-operator} is precisely the mismatch between the two:
it vanishes exactly when $\mathcal{D}$ reproduces the model's own endogenous
surface dynamics, and is otherwise the signed amount by which the hedging
Greeks fail the model's own equation. The pair $(P,\,\partial P_{\mathrm{model}})$
solves \eqref{pe:fk}; the pair $(P,\,\partial P_{\mathrm{hyb}})$ does not ---
same price, sensitivities from a different world.
\trace{Correctness; VM-4, VM-5; Vol~I \S2.4.2, Eq.~(2.21), \S3.6, Eq.~(3.34).}

%% ---------------------------------------------------------------------------
\pelabel{PE-3}\textbf{PE-3. Product instantaneous volatility: the coefficient
that restores the equation, and why it need not be smooth.} The discrepancy
\eqref{pe:residual-operator} can be absorbed into a single coefficient. Since
$\smod$ enters \eqref{pe:fk} only through $\tfrac12\smod^2S^2\partial_{SS}P$,
one asks for the diffusion coefficient that makes a PDE \emph{of the same form}
hold with the model theta and the hybrid delta and gamma.

\textbf{Definition (product instantaneous variance and volatility).} At a point
$(S,t)$ with $\Ghyb(S,t)\neq 0$, the \emph{product instantaneous variance}
$\sprod^2(S,t)$ is the sign-carrying number solving
\begin{equation}\label{pe:sprod-pde}
  \partial_t P + \tfrac12\,\sprod^2 S^2\,\Ghyb + (r-q)S\,\Dhyb - rP = 0 ,
\end{equation}
that is, in closed form,
\begin{equation}\label{pe:sprod}
  \boxed{\;
  \sprod^2 \;=\; -\,\frac{2\bigl[\,\Theta + (r-q)S\,\Dhyb - rP\,\bigr]}{S^2\,\Ghyb}
  \;,}\qquad \Theta = \partial_t P .
\end{equation}
It is the first-class object here --- real and total wherever $\Ghyb\neq0$, its
sign a diagnostic in its own right; the \emph{product instantaneous volatility}
$\sprod=\sqrt{\sprod^2}$ is its nonnegative root where $\sprod^2\ge0$. Under
$r=q=0$ this is $\sprod^2 = -2\Theta/(S^2\Ghyb)$; with $q=0$ it matches
the normative form $\sprod^2 = -2(\theta + rS\Dhyb - rP)/(S^2\Ghyb)$.

\textbf{Well-posedness (stated, not waved).} Equation \eqref{pe:sprod} defines
$\sprod^2$ wherever $\Ghyb\neq0$. Writing $N := -2[\Theta+(r-q)S\Dhyb-rP]$ and
$M := S^2\Ghyb$, the value $\sprod=\sqrt{N/M}$ is real and nonnegative iff $N$
and $M$ share sign; where $M=0$ the ratio has a pole and $\sprod$ is undefined;
where $N/M<0$ no real $\sprod$ exists --- the FK form cannot be restored with a
real coefficient; and where $\mathcal{D}$ carries $\Sigma(\cdot\,;S)$ out of the
admissible set $\mathcal{A}$ (PE-1, A5), the model price of a surface-dependent
product is itself undefined, and $\sprod$ inherits that gap. These are
structural facts about the hybrid pair, recorded, not smoothed.

\textbf{Its nature, stated with total honesty.}
\emph{(i) Reverse-engineered, not an SDE coefficient.} $\sprod$ is defined by
algebraically inverting \eqref{pe:sprod-pde} given $(\Theta,\Dhyb,\Ghyb)$; it is
not the diffusion coefficient of any stochastic differential equation for $S$.
In particular $\sprod \neq \smod$: it is neither the constant $\sigma$ nor
$\sigma_{\mathrm{loc}}$.
\emph{(ii) Specific to the product and to the rule.} Through $\Dhyb,\Ghyb$ the
right-hand side of \eqref{pe:sprod} depends on the product $\Phi$ (via its
surface sensitivities) and on $\mathcal{D}$ (via $\partial_S\Sigma,
\partial_{SS}\Sigma$). Two products at one $(S,t)$ under one rule have different
$\sprod$; one product under two rules (say sticky-strike versus sticky-delta)
has different $\sprod$.
\emph{(iii) Fitted, hence not required to be smooth --- and here is why.}
$\sprod^2$ is the ratio $N/M$ of two independently constructed fields; it solves
no differential equation that would impose regularity, so its only regularity is
that inherited from the quotient. The denominator $M=S^2\Ghyb$ can vanish and
change sign --- Vol~I exhibits $\Gadj$ crossing zero as Vanna does
(Vol~I, \S3.9.3) --- while the numerator $N$ need not, so
$\sprod^2$ carries poles and sign changes there. Even where finite, $N$ and $M$
inherit $\partial_S\Sigma$ and $\partial_{SS}\Sigma$ from the declared rule
$\mathcal{D}$, itself only a declared map --- admissibly kinked (a
sticky rule piecewise in moneyness) --- so nothing forces $\sprod$ to be even
continuous. Contrast the Dupire local volatility, which is a smooth functional
of a smooth arbitrage-free surface (Vol~I, \S2.5): $\sigma_{\mathrm{loc}}$ is
\emph{modelled}; $\sprod$ is \emph{fitted}, and the two regularities are not the
same.
\trace{Correctness; VM-4, VM-6; Vol~I \S2.5, \S3.9.}

%% ---------------------------------------------------------------------------
\pelabel{PE-4}\textbf{PE-4. Carry: theta is model-determined; the theta--gamma
balance closes only against the product instantaneous volatility.}

\textbf{Claim 1 (theta is invariant to the surface-vs-spot dynamic).} The model
theta $\Theta=\partial_t P$ does not depend on the choice of $\mathcal{D}$.

\textbf{Derivation.} $\Theta$ is the partial of the model price in calendar time
at \emph{fixed} spot. By (A5) $\mathcal{D}$ is a spot-dynamic: it prescribes the
surface response $\partial_S\Sigma$ and enters the reprice map only through
$S$-derivatives, where the chain rule multiplies the surface terms
$D_\Sigma P[\cdot]$ by $\partial_S\Sigma$ --- see
\eqref{pe:dhyb}--\eqref{pe:ghyb}. The calendar partial at fixed $S$ contains no
$\partial_S\Sigma$ factor; hence $\Theta$ is one and the same across the whole
admissible family of surface-move maps, while $\Dhyb,\Ghyb$ vary with
$\mathcal{D}$ through $\delta_\Delta,\delta_\Gamma$. This is the asymmetry the
normative source states: the dynamics of the implied surface with respect to
spot affect the hybrid delta and gamma and leave theta untouched. $\square$

\textbf{Scope (honest caveat).} The invariance holds because $\mathcal{D}$ is
declared as a spot-dynamic. A rule that additionally prescribed an autonomous
calendar roll $\partial_t\Sigma$ independent of the model would give a
``hybrid theta'' carrying $D_\Sigma P[\partial_t\Sigma]$; PE-1(A5) fixes
$\mathcal{D}$ as a spot-dynamic precisely so that theta is the model's alone.

\textbf{Claim 2 (the theta--gamma balance is readable only at $\sprod$).}
Rearranging \eqref{pe:sprod-pde},
\begin{equation}\label{pe:carry-identity}
  \Theta + (r-q)S\,\Dhyb - rP \;=\; -\tfrac12\,\sprod^2 S^2\,\Ghyb ,
\end{equation}
and under the zero-rate convention ($r=q=0$, $F=S$, flagged),
$\Theta = -\tfrac12\,\sprod^2 S^2\,\Ghyb$. The hedger carries $\Ghyb$; the daily
breakeven move --- the $|\delta S|$ at which gamma gain offsets a day's decay ---
follows from $\tfrac12\Ghyb(\delta S_{\mathrm{BE}})^2 = |\Theta|\,\dd t$ and
\eqref{pe:carry-identity} as
\begin{equation}\label{pe:breakeven}
  \delta S_{\mathrm{BE}} = \sprod\,S\,\sqrt{\dd t}\,,
\end{equation}
the same form as the corpus breakeven $\sigma S/\sqrt{252}$ (Vol~I, \S1.4,
Eq.~(1.13)) but with $\sprod$ in place of the implied $\sigma$ --- and it
presumes the real, long-carry regime $\sprod^2>0$; where PE-3 gives
$\sprod^2\le0$ (a pole, or a sign flip) there is no real breakeven. So $\sprod$
is, by construction, the one volatility at which the model theta and the
hedger's gamma reconcile: a volatility $\sigma$ leaves in
\eqref{pe:carry-identity} the residual $\tfrac12(\sigma^2-\sprod^2)S^2\Ghyb$,
nonzero for every $\sigma\neq\sprod$ and equal to the Feynman--Kac discrepancy
$\mathcal{L}_{\mathrm{hyb}}$ of \eqref{pe:residual-operator} exactly at the
model's own $\smod$.

\textbf{Linkage to Part~A (the statement threaded there).} The
profit-and-loss explain certificate's carry line reconciles the model theta
against the gamma the book actually carries --- the hybrid gamma. That
reconciliation holds at exactly one volatility, $\sprod$. A certificate that
reconciles at a volatility other than $\sprod$ carries a nonzero residual into
$R$ --- the model's $\smod$ leaves exactly $\mathcal{L}_{\mathrm{hyb}}\,\dd t$ of
\eqref{pe:residual-operator}, a general $\sigma$ leaves
$\tfrac12(\sigma^2-\sprod^2)S^2\Ghyb\,\dd t$ --- which is exactly the unexplained
profit-and-loss VM-6 forbids to leave silently unattributed and VM-7 calls a
broken chain. Hence the carry line carries a declared, recorded term naming
which volatility reconciles it; VM-4 and VM-6 carry that thread.
\trace{Correctness, Auditability; VM-4, VM-6, VM-7; Vol~I \S1.4, Eq.~(1.13).}

%% ---------------------------------------------------------------------------
\pelabel{PE-5}\textbf{PE-5. The corpus illustration: the total gamma under an
implied-volatility dynamic drives the product instantaneous volatility,
vanillas included.} Take a single name, the constant-vol model priced at the
option's implied volatility, and the corpus surface-move map. All under the
zero-rate convention $F=S$, flagged.

\textbf{The declared dynamic, explicitly.} Write the implied vol at the option's
fixed strike through the corpus decomposition (Vol~I, \S2.1.1, Eq.~(2.1))
$\sigma(m,\lF,T) = \satm(\lF,T) + \ssm(m,T)$, with $m=\ln(K/F)$,
$\lF=\ln(F/\Fref)$, skew $s(m)=\partial_m\ssm$, curvature
$c(m)=\partial_m^2\ssm$, and ATM slope $\volslope=\partial_{\lF}\satm$. The map
$\mathcal{D}$ is carried entirely by $(\volslope,\,s,\,c)$: routing through
$F(S)$ gives (Vol~I, \S2.4.2, Eqs.~(2.19)--(2.20))
\begin{equation}\label{pe:dsig}
  \frac{\dd\sigma}{\dd S} = \frac{\volslope - s(m)}{S},
  \qquad
  \frac{\dd^2\sigma}{\dd S^2}
  = \frac{\bigl(\partial_{\lF}^2\satm - \volslope\bigr) + \bigl(c(m)+s(m)\bigr)}{S^2}.
\end{equation}

\textbf{The hybrid gamma, every cross term explicit.} The price is
$C(S,\sigma(S))$ with $\sigma(S)$ read off the surface under $\mathcal{D}$. Its
second total spot derivative is the total (smile-adjusted) gamma of Vol~I, \S2.4.2,
Eq.~(2.21):
\begin{equation}\label{pe:gadj}
  \Ghyb = \Gadj
  = \Gamma + 2\,\Vanna\,\frac{\dd\sigma}{\dd S}
    + \Volga\left(\frac{\dd\sigma}{\dd S}\right)^{\!2}
    + \nu\,\frac{\dd^2\sigma}{\dd S^2},
\end{equation}
with $\Gamma=\partial_{SS}C$, $\Vanna=\partial_S\partial_\sigma C$,
$\Volga=\partial_\sigma^2 C$, $\nu=\partial_\sigma C$ the frozen-vol
Black--Scholes Greeks. This is exactly the $\delta_\Gamma$ of
\eqref{pe:ghyb} in the single-implied-vol realization: the model gamma
$\Gamma$, plus the vanna and volga terms in the surface velocity
$\dd\sigma/\dd S$, plus vega times the surface acceleration $\dd^2\sigma/\dd
S^2$. It is confirmed to be the object that enters the P\&L identity in
Vol~I, \S3.6, Eq.~(3.34).

\textbf{The product instantaneous volatility, exhibited.} The constant-vol model
theta is $\Theta = \partial_t C = -\tfrac12\sigma^2 S^2\Gamma$ (Feynman--Kac
\eqref{pe:fk} at $r=q=0$ with $\smod=\sigma$, $\partial_{SS}C=\Gamma$). Insert
into \eqref{pe:sprod} with $\Ghyb=\Gadj$:
\begin{equation}\label{pe:sprod-vanilla}
  \sprod^2 = \frac{-2\Theta}{S^2\Gadj}
           = \sigma^2\,\frac{\Gamma}{\Gadj},
  \qquad\text{i.e.}\qquad
  \sprod = \sigma\,\sqrt{\frac{\Gamma}{\Gadj}} .
\end{equation}
The model theta fixes the numerator and is invariant to $\mathcal{D}$ (PE-4);
the entire dependence on the assumed dynamic sits in the ratio
$\Gamma/\Gadj$, and by \eqref{pe:gadj} that ratio departs from $1$ for
\emph{any} product carrying vega, vanna, or volga once the surface moves with
spot ($\volslope\neq0$ or $s(m)\neq0$). A vanilla carries all three; therefore
$\sprod\neq\sigma$ for a vanilla under any nontrivial surface dynamic --- the
normative source's ``vanilla options included'' is exactly \eqref{pe:sprod-vanilla}.
The clean ratio is specific to the constant-vol model: for a local-vol model
$\sprod$ is the general form \eqref{pe:sprod} with
$\Theta=-\tfrac12\sigma_{\mathrm{loc}}^2 S^2\,\partial_{SS}P$ and
$\partial_{SS}P$ the local-vol price's own second partial --- the
$\sqrt{\Gamma/\Gadj}$ collapse does not survive.

\textbf{A number.} Take Part~A's base call ($\sigma=30\%$, $\Gamma=0.0265$, vega
$\approx0.20$ per vol point, $S=100$, $T=0.25$) under a downward skew of one vol
point per one-percent up-move ($\dd\sigma/\dd S\approx-0.01$). Its vanna is
$\approx0.10$, so the surface-velocity term
$2\,\Vanna\,(\dd\sigma/\dd S)\approx-0.002$ lowers the effective gamma to
$\Gadj\approx0.0245$, and
$\sprod=\sigma\sqrt{\Gamma/\Gadj}\approx30\%\times\sqrt{0.0265/0.0245}\approx31\%$.
The breakeven volatility the delta-hedged holder actually runs is about $31\%$,
not the $30\%$ quoted --- a vol point of carry the implied number never showed.

\textbf{Non-smoothness, exhibited.} Vol~I, \S3.9.3, exhibits $\Gadj$ running from a
positive smile adjustment at one spot to a negative one a finite move away, as
Vanna crosses zero. Where $\Gadj\to0$, \eqref{pe:sprod-vanilla}
gives $\sprod\to\infty$ (a pole); where $\Gadj$ takes the sign opposite to
$\Gamma$, $\sprod^2<0$ and no real $\sprod$ exists. So even for a vanilla,
$\sprod$ is neither globally smooth nor globally real --- the abstract statement
of PE-3(iii) made concrete.

\textbf{The multi-name instance.} When the book's value depends on an implied
correlation re-marked by a declared convention, the surface-response term
$\delta_\Gamma$ of \eqref{pe:ghyb} acquires a correlation channel: the
\emph{shadow gamma} of Vol~II, \S2.5, Eq.~(2.23) --- under the equal-vol,
locally-constant-$\volslope$ hypothesis of its Prop.~2.27 ---
\[
\begin{aligned}
  \mathrm{ShadowGamma}
  &= \frac{\volslope^2}{S}\!\left[\partial_\rho^2 V\,R^2
      + \partial_\rho V\,\partial_\sigma^2\rho\big|_\xi\right]
   + \volslope R\!\left(\partial^2_{S\rho}V
      + \frac{\volslope}{S}\partial^2_{\sigma\rho}V\right),\\
  &\qquad\qquad R=\partial_\sigma\rho\big|_\xi,
\end{aligned}
\]
under the sticky-CVC convention ($\dd\xi=0$), and vanishing identically under
sticky-$\rho$ ($\dd\rho=0$). It enters $\Ghyb$ on its raw footing
$\Gamma^{\mathrm{sh}}$ (Vol~II, \S2.5, Eq.~(2.21)), so that here too the choice
of surface-move convention drives $\sprod$ through \eqref{pe:sprod} --- the same
structure as the single-name case, with $\mathrm{Cega}$ and the correlation
dynamics in place of vega and the vol dynamics.
\trace{Correctness; VM-4, VM-5; Vol~I \S2.4.2, Eqs.~(2.19)--(2.21), \S3.9.3;
  Vol~II \S2.5, Eqs.~(2.21),(2.23), Prop.~2.27.}

%% ---------------------------------------------------------------------------
\pelabel{PE-6}\textbf{PE-6. The product instantaneous variance is a first-class
derived quantity: computed on demand per product and per valuation over the
chain.} The canonical object is the \emph{product instantaneous variance}
$\sprod^2$ of \eqref{pe:sprod} --- the sign-carrying effective variance, real and
total wherever $\Ghyb\neq0$ --- with the product instantaneous volatility
$\sprod=\sqrt{\sprod^2}$ its nonnegative root where $\sprod^2\ge0$. It is a
\emph{projection} of the recorded Greeks: the deterministic algebraic function
\eqref{pe:sprod} of $\Theta,\Dhyb,\Ghyb,P$ and the market state, so it
reproduces \emph{unconditionally} from the record with no model re-run ---
stronger than a mark, whose number needs its retained model (C-14.15). Being a
projection it is computed when needed and never stored (C-4.11, MD-6): for every
product and every valuation its value, its sign, its poles, and its regularity
are monitored by \emph{recomputation over the chain's recorded Greeks}, which is
authoritative, so it cannot drift (VM-8). Where the object is pathological the
pathology is what the recomputation returns: at a pole ($\Ghyb=0$) $\sprod^2$ is
undefined and the fact is the pole itself; where $\sprod^2<0$ the fact is that no
real volatility restores the Feynman--Kac form --- the sign is the diagnostic,
not an error to be smoothed. Its evolution across links is thereby auditable and
reproducible under challenge, an internal diagnostic feeding the certificate of
a dispute-ready mark (VM-8), not itself a counterparty mark. \emph{Why the
mandate:} a pole or a sign flip in $\sprod^2$ marks \emph{degenerate carry} ---
the spot at which the carry line's theta-gamma reconciliation breaks down, the
carry-line analogue of the $O(1)$ residual VM-4 books at a digital. This is the
single monitoring mandate of Part~B; it prescribes nothing further.
\emph{Parked:} the mandate that $\sprod^2$ additionally be \emph{stored} per
valuation on the chain is parked against C-4.11 and MD-6 --- a projection stores
nothing --- and stands in the open-problems index for the owner's ruling;
nothing above depends on storage.
\trace{Auditability, Reproducibility; VM-1, VM-3, VM-6, VM-8, C-4.11, C-14.15, MD-6, MD-15.}

\section{Open Problems (Parked Items)}

The framework parks a genuine conflict rather than fudge it: it is recorded here
with the exact clause it would replace, the exact amendment text, and the
justification, and the decision is the owner's alone. This index holds the
document's first entry.

\textbf{PARK-1 --- storing the product instantaneous variance versus C-4.11.}

\emph{Conflict.} Part~B mandates that $\sprod^2$ (PE-6) be \emph{stored} per
product and per valuation on the chain. This conflicts with the Constitution's
C-4.11 and with MD-6: $\sprod^2$ is a \emph{projection} --- an algebraic
recomputation over recorded Greeks, needing no model run (PE-6) --- and a
projection stores nothing. It would be the architecture's first stored
ledger-computable output. The in-force text of PE-6 therefore runs the
conforming reading --- computed on demand, monitored by recomputation, never
stored --- and only the storage half of the mandate is parked.

\emph{Exact clause it would replace (C-4.11, final sentence).} ``Anything
computable from the coordinates is a projection, computed when needed and never
stored.''

\emph{Exact proposed amendment text.} ``Anything computable from the coordinates
is a projection, computed when needed; recomputation from the record is always
authoritative. A projection's value may be recorded as a verifiable field of a
re-entered observation or a certificate --- a \textbf{materialised projection}
--- only where (i) recomputation from the record remains authoritative and always
available, (ii) the recorded value carries complete lineage to the inputs it was
computed from, and (iii) any correction to those inputs flags the recorded value
stale. So recorded, it is never a second source that can silently diverge, and
recomputation is the standing check. Absent all three, a projection is never
stored.'' Mirrored in MD-6: ``stores nothing'' becomes ``stores nothing, save as
a materialised projection under C-4.11's three conditions.''

\emph{Justification.} A recomputable, lineage-bearing, flagged-stale field cannot
silently diverge, so the reconciliation-failure mode the architecture exists to
remove (C-1.2, C-1.3) stays closed and a materialised $\sprod^2$ series would be
operationally safe. Safety is not textual permission: C-4.11 is categorical, and
amending it is the owner's alone.

\emph{Alternative resolution, keeping C-4.11 literal.} The storage mandate's
``stored'' becomes ``computed on demand and monitored,'' $\sprod^2$ a projection
never stored --- exactly the conforming reading PE-6 already runs; under it no
amendment is needed and the park closes.

\end{document}
